\scanpage{9}
% L1–7
L'image \inserted{essentielle} de $\alpha$ est la sous-catégorie \inserted{str} pleine de $\prod_i\mathrm{Hom\,top}(X,E_i)$ formée des systèmes $(f_i\colon X\longrightarrow E_i)$ tels que pour tt $i$, $f_i$ se factorise par $E'_i$ ; \unsure{ou|or} \unsure{encore|alors}, \deleted{\unsure{donc $f$}} par $\beta$, elle correspond à la ss-cat. str. pleine de $\underline{\mathrm{Hom\,top}}\,(X,E)$ formée des $f\colon X\longrightarrow E$ tels que \inserted{$\forall i$,} $\mathrm{pr}_i\circ f$ se factorise par $E'_i$, \uline{i.e.} tels que \inserted{$\forall i$,} $f$ se factorise par $F_i$, i.e. tels que $X$ se factorise par $F$. Donc

% L8
\[
\prod_i\underline{\mathrm{Hom\,top}}\,(X,E'_i)\cong\underline{\mathrm{Hom\,top}}\,(X,F).
\]

% L9–10
\uline{Corollaire de} a) et b) : Le 2-produit d'une famille de topos existe tjrs.

% L11–12
\deleted{Pour avoir la $\varprojlim_{\mathrm{top}}$ d'un topos fibré, il faut}

% L12–15-left
\diagram{\begin{tikzcd}[ampersand replacement=\&,row sep=small]
\& \deleted{\mathscr F_j} \& \deleted{\mathscr F_i}\arrow[l]\\
\deleted{\mathscr F}\arrow[d,dash] \& \deleted{\mathscr F_j}\arrow[r,"{\deleted{f^*}}"] \& \deleted{\mathscr F_i}\\
\deleted{I} \& \deleted{j} \& \deleted{i}\arrow[l,"{\deleted{f}}"]
\end{tikzcd}}

% L13–16-centre
\[
\begin{aligned}
&\deleted{\varprojlim_i\underline{\mathrm{Hom\,top}}\,(X,\mathscr F_i)}\\[-2pt]
&\deleted{\rotatebox{90}{\ensuremath{\simeq}}}\\[-2pt]
&\deleted{\varprojlim_i\underline{\mathrm{Hom}}^{\unsure{!|'}}(\mathscr F_i,X)}\\[-2pt]
&\deleted{\rotatebox{90}{\ensuremath{\simeq}}}
\end{aligned}
\]

% L15–17-right
\diagram{\begin{tikzcd}[ampersand replacement=\&,row sep=tiny,column sep=small]
\deleted{x}\arrow[r,bend left,dash] \& \deleted{y}\arrow[r,bend left,dash] \& \deleted{z}\arrow[r,"\deleted{\sim}"] \& {}\\
\& \deleted{\unsure{f^*(x)}}\arrow[u,phantom,"\simeq"] \& \deleted{g^*f^*(x)}\arrow[r,"\deleted{\sim}"']\arrow[u,phantom,"\unsure{1}"] \& \deleted{g^*(y)}\\
{}\arrow[r,"{\deleted{f}}"'] \& {}\arrow[l] \& {}\arrow[l,"{\deleted{g}}"'] \& {}\\
{}\arrow[r,dash] \& {} \& {} \& {}
\end{tikzcd}}

% L18
\deleted{$(E_i)_{i\in I}$ pseudo syst projectif de topos}

% L18–21
\diagram{\begin{tikzcd}[ampersand replacement=\&,row sep=small,column sep=small]
\& \deleted{\mathscr F_1}\arrow[r,dash,shift left]\arrow[r,dash,shift right] \& \deleted{\mathscr F_2}\arrow[r,dash,shift left=2]\arrow[r,dash]\arrow[r,dash,shift right=2] \& \deleted{\mathscr F_3}\\
\deleted{\underset{i}{\underline{\lim}}_{\mathrm{top}} E_i}\arrow[r] \& \deleted{\prod_{i\in\mathrm{Ob}\, I}E_i}\arrow[r,shift left]\arrow[r,shift right]\arrow[d,"{\deleted{\mathrm{pr}_{s(\alpha)}}}"']\arrow[dr,"{\deleted{\mathrm{pr}_{\unsure{s(\beta)|b(\alpha)}}}}"] \& \deleted{\prod_{\alpha\in\mathrm{Fl}_1 I}E_{b(\alpha)}}\arrow[r,shift left=2]\arrow[r]\arrow[r,shift right=2] \& \deleted{\prod_{(\alpha,\beta)\in\mathrm{Fl}_2(I}E_{b(\beta)}}\\
\& \deleted{E_{s(\alpha)}}\arrow[r] \& \deleted{E_{b(\alpha)}}\arrow[u,phantom,"{\deleted{\text{(but}}}"] \&
\end{tikzcd}}

% foot
% archive sheet number: 8
